\documentclass[journal=jpclcd, manuscript=letter, layout=twocolumn, doi=true]{achemso}

\usepackage[utf8]{inputenc}
\usepackage{amsmath}
\usepackage{amssymb}
\usepackage{graphicx}
\usepackage{physics}
\usepackage{booktabs}
\usepackage{xcolor}
\usepackage{subcaption}
\usepackage{algorithm}
\usepackage{algpseudocode}
\usepackage[hidelinks]{hyperref}
\newcommand{\red}{\color[rgb]{0.8,0,0}}
\newcommand{\green}{\color[rgb]{0.0,0.6,0.0}}
\newcommand{\blue}{\color[rgb]{0.0,0.0,0.6}}
\usepackage[hidelinks]{hyperref}

\title{Non-Equilibrium Dynamics of the Time-Dependent Excitonic Coupling in Fluorescent Protein Dimers}

%\title{Interplay of Time-Dependent Excitonic Coupling and Decoherence in Protein Dimers}

\author{Robson Christie}
\affiliation{School of Mathematics and Physics, University of Portsmouth, PO1 3FX, United Kingdom}

\author{Cerys Murray}
\affiliation{School of Mathematics and Physics, University of Surrey, Guildford GU2 7XH, UK}

\author{Youngchan Kim}
%\email{youngchan.kim@surrey.ac.uk}
\affiliation{Leverhulme Quantum Biology Doctoral Training Centre, School of Biosciences, Advanced Technology Institute, University of Surrey, Guildford GU2 7XH, UK}

\author{Jaewoo Joo}
\email{* jaewoo.joo@port.ac.uk}
\affiliation{School of Mathematics and Physics, University of Portsmouth, PO1 3FX, United Kingdom}

\begin{document}

\maketitle

% Move tocentry here to bypass layout-induced frame distortion
\begin{tocentry}
\includegraphics[width=1.0\textwidth]{dimer_ray_clean_1024.png}
\end{tocentry}

\begin{abstract}
\textbf{Corrected version (v2).} \emph{Version 1 of this preprint reported
$J = 74.38~\mathrm{cm^{-1}}$ and a point-dipole comparison of
$13.31~\mathrm{cm^{-1}}$, and attributed the resulting $5.6\times$ ratio to
near-field multipolar effects. That ratio was an artefact of a unit conversion
in the Coulomb sum, and the two numbers were not computed on matched geometries.
The corrected values and the conclusions that replace them are given below; a
full claim-by-claim correction accompanies this version. Readers who have cited
the $5.6\times$ near-field result should use this version instead.}

We quantify the excitonic coupling in the homodimer of dimeric Venus fluorescent
protein using a quantum-classical hybrid workflow. Employing a transition-density
coupling (TDC) formalism over a 1000-frame molecular dynamics ensemble, we obtain
$J = 32.82 \pm 1.55~\mathrm{cm^{-1}}$ under optical-limit screening
($\varepsilon_{\mathrm{opt}} = 1.77$) at a chromophore centroid separation of
$24.69 \pm 0.32~\mathrm{\AA}$. Evaluated on the same frames and the same
transition density, the far-field point-dipole approximation (PDA) gives
$27.64 \pm 1.28~\mathrm{cm^{-1}}$, a ratio of $1.187 \pm 0.011$. Near-field
multipolar corrections are therefore a $\sim 19\%$ effect at this separation, not
a dominant one: the PDA is adequate for fluorescent-protein dimers at biological
separations, and cheap estimates are sufficient. We further find the dimer to be
detuning-dominated rather than delocalised, with an ensemble site-energy gap
$\langle|\Delta|\rangle = 576~\mathrm{cm^{-1}}$ greatly exceeding
$2|J| = 65.6~\mathrm{cm^{-1}}$, giving localised eigenstates with a mean mixing
angle of $0.18$. This accounts for the reported spectroscopy without invoking a
protected delocalised exciton, and we show that the circular-dichroism couplet
separation is not an estimator of the exciton gap in either direction, being a
non-monotonic and lineshape-dominated function of it. Comparison against an
independent two-photon polarisation measurement agrees to within the dielectric
screening convention, which we identify as the dominant remaining systematic ---
larger than the near-field correction originally reported. Finally, we argue that
a separation of timescales reconciles robust excitonic signatures with a strongly
decohering biological environment, irrespective of the degree of environmental
noise suppression, and employ stochastic simulations to model the post-absorption
evolution into incoherent hopping between localised chromophore states.
\end{abstract}

\section{Introduction}
Excitation energy transfer in biological systems is conventionally understood through the competition between intermolecular excitonic coupling ($J$) and environmental vibrational dephasing \cite{forster1965delocalized, may2004charge}. When environmental fluctuations dominate, phase relationships randomise rapidly, restricting the system to the very weak-coupling regime characterised by incoherent Förster hopping \cite{forster1948intermolecular}. Conversely, the intermediate weak regime permits the formation of coherently delocalised excitonic eigenstates, manifesting spectroscopically as Davydov splitting \cite{davydov1971}, while the strong coupling regime gives rise to more pronounced modifications of the absorption spectrum \cite{clegg2006history}. Because physiological environments are highly dynamic and structurally disordered, biophysics has historically assumed that ambient thermal noise precludes robust excitonic delocalisation, confining such systems to the very weak coupling regime \cite{clegg2006history, nelson2018role}.

Homodimers of the yellow fluorescent protein variant dimeric Venus challenge this assumption. Operating at room temperature, these dimers exhibit distinct Davydov splitting in circular dichroism (CD) spectra, and photon antibunching statistics confirm that the interacting chromophores act as a collective quantum entity upon photon absorption \cite{kim2019venusa206, freed2025effect}. These signatures indicate an interaction strength firmly within the intermediate excitonic regime, seemingly at odds with the biological environment inducing strong decoherence.

To explain this phenomenon, it has been hypothesised that the fluorescent protein $\beta$-barrel scaffold plays a central role in sustaining inter-chromophore excitonic coupling by providing a structured dielectric shield that attenuates thermal fluctuations from the surrounding environment, thereby potentially extending excitonic decoherence times \cite{kim2019venusa206}. While this protective shielding offers a compelling mechanism, we propose that the robust excitonic coupling can be fundamentally reconciled by separating the physical timescales of exciton formation and subsequent dephasing, irrespective of the exact degree of environmental noise suppression. Photoexcitation acts as a collective, near-instantaneous vertical transition directly into a delocalised bright eigenstate, initially protected by the scaffold's optical-limit dielectric screening. The Davydov splitting is an inherent property of the energy spectrum imprinted at this precise moment of absorption. As demonstrated by the stochastic trajectories in this work (see {\bf Figure \ref{fig:ExcitonDynamics}}), the system subsequently undergoes strong environmental dephasing and rapid relaxation into a localised state. However, because this decoherence and the associated dielectric screening occur on sub-picosecond to picosecond timescales, they are significantly slower than the absorption event. Consequently, strong excitonic coupling and its spectroscopic signatures are observed despite the presence of strong, subsequent decoherence.

Predicting the magnitude of the coupling that drives this initial delocalisation is often held to require going beyond standard point-dipole interaction models, since at short intermolecular separations near-field multipolar effects contribute appreciably to the coupling strength $J$ \cite{scholes2003, kitohnishioka2017forster}. Capturing them requires integrating the full three-dimensional (3D) spatial distribution of the electron transition density, which maps the shape of the transition's probability amplitude onto the 3D molecular structure. A question this work settles quantitatively is whether that expense is warranted at the separations relevant to fluorescent-protein dimers, where the chromophores are held apart by their $\beta$-barrel scaffolds; we find that it is not.

Furthermore, because exciton formation precedes the macroscopic nuclear reorganisation of the solvent (Debye relaxation $\tau_D \approx 8.3$\,ps), the initial Coulombic coupling is not damped by the static permittivity of water ($\varepsilon_{\mathrm{eq}} \approx 78$). Instead, it is screened only by the fast electronic polarisation of the environment, corresponding to the optical dielectric limit ($\varepsilon_{\mathrm{opt}} \approx 1.78$) \cite{mennucci2012, scholes2003}.

In this work, we present a quantum-classical hybrid analysis of non-equilibrium dynamics of the time-dependent excitonic coupling between dimeric Venus molecules. By combining a transition-density coupling (TDC) formalism with a non-equilibrium dielectric framework, we quantify the near-field multipolar contribution to the excitonic coupling and find it to be a $\sim 19\%$ correction at this separation. We show that the resulting coupling is small compared with the ensemble site-energy detuning, so that the eigenstates are localised rather than delocalised, and that the dominant systematics in any comparison with experiment are the dielectric screening convention and the monomer transition dipole rather than the multipolar correction itself.
%%%%%%%%%%%%%%%%%%%%%%%%%%%%%%%%%%%%%%%%%%%%%%%%
\section{Theoretical Framework}

\begin{figure*}[t]
  \centering
  \begin{subfigure}[b]{0.32\textwidth}
    \centering
    \includegraphics[width=\linewidth]{Fig_Bloch_Grid.pdf}
    \caption{}
    \label{fig:bloch}
  \end{subfigure}\hfill
  \begin{subfigure}[b]{0.32\textwidth}
    \centering
    \includegraphics[width=\linewidth]{Fig_SSE_Site.pdf}
    \caption{}
    \label{fig:sse_site}
  \end{subfigure}\hfill
  \begin{subfigure}[b]{0.32\textwidth}
    \centering
    \includegraphics[width=\linewidth]{Fig_SSE_Adiabatic.pdf}
    \caption{}
    \label{fig:sse_adiabatic}
  \end{subfigure}

  \caption{ 
  Simulation of excitonic dynamics in the Venus dimer with the local site states $\ket{1}\equiv\ket{e_L g_R}$ and $\ket{2}\equiv\ket{g_L e_R}$: (a) Unified Bloch sphere showing the stochastic pure-state trajectory (red) and ensemble-averaged mixed-state path (blue); (b) Site-basis populations and coherence for a single SSE trajectory; (c) Adiabatic populations of bright and dark states $\ket{\pm}$ for a single SSE trajectory.}
  \label{fig:ExcitonDynamics}
\end{figure*}

To describe the excitonic interactions within the dimer, we restrict our model to the single-exciton manifold, defining the local site basis as $\ket{1}\equiv\ket{e_L \, g_R}$ and $\ket{2}\equiv\ket{g_L \, e_R}$. This notation explicitly enforces the conservation of energy in a single-photon absorption subspace, ensuring that only a single excitation is shared between the two chromophores. The beta-barrel of Venus has a diameter of $\sim 22$~\AA\ and a closest separation distance of 15~\AA\ in the crystal structure of the dimeric Venus fluorescent protein homodimer~\cite{kim2019venusa206}, effectively precluding significant short-range Dexter exchange interactions~\cite{skourtis2016dexter}. Consequently, the electronic coupling $J$ is accurately described as the Coulomb interaction between the transition densities of the individual monomers \cite{hsu2009}. The electron transition density, $\rho^{g\to e}_m(\mathbf r)$, is obtained by integrating the product of the ground and excited state wavefunctions over the coordinates of $N-1$ electrons in a monomer:
\begin{multline} \label{eq:transition_density}
    \rho^{g\to e}_m(\mathbf{r}) =  
    N \int \dots \int d\mathbf{r}_2 \dots d\mathbf{r}_N \\ \Psi_e^*(\mathbf{r}, \mathbf{r}_2, \dots, \mathbf{r}_N) \, \Psi_g(\mathbf{r}, \mathbf{r}_2, \dots, \mathbf{r}_N)
\end{multline}
where $\Psi_g$ and $\Psi_e$ are the many-body electronic wavefunctions for the ground and excited states, respectively ($m=L,\,R$). 
Crucially, while the total electron density is a positive-definite observable representing the probability of finding an electron, the electron transition density implies an off-diagonal matrix element. Since $\int \rho^{g\to e}(\mathbf{r}) d\mathbf{r} = 0$ due to the orthogonality of the ground and excited states, its spatial distribution, characterised by positive and negative phases, indicates the transition's multipolar character. For example, weighting the transition density by the position vector yields the transition dipole moment ($\boldsymbol{\mu} = \int \mathbf{r} \rho^{g \to e}(\mathbf{r}) d\mathbf{r}$), which dictates the far-field optical response. In order to compute the full TDC, the 3D spatial integral is given by 
\begin{equation}
    J_{\mathrm{TDC}} \approx \frac{1}{4\pi\varepsilon(t)}\iint \frac{\left(\rho^{g\to e}_L(\mathbf r)\right)^*\,\rho^{g\to e}_R(\mathbf r')}{|\mathbf r-\mathbf r'|} \,d\mathbf r\,d\mathbf r'
    \label{eq:J_integral}
\end{equation}
for time-dependent permittivity $\varepsilon(t)$.
Formally, the electron transition density can be complex, and the conjugate ensures a real coupling strength. In our implementation, the space is discretised over a grid as $J \approx\frac{1}{4\pi\varepsilon(t)}\sum_{i\in L}\sum_{j\in R}\frac{q_i q_j}{|\mathbf r_i-\mathbf r_j|}$, where $q_i$ and $q_j$ represent the discrete transition atomic charges. For comparison, we also evaluate the far-field PDA, which collapses these densities into two idealised point dipoles $\boldsymbol{\mu}_m$ separated by vector $\mathbf{R}$ given by
\begin{equation} \label{eq:J_PD}
    J_{\mathrm{PDA}} = \frac{1}{4\pi\varepsilon(t)} \left( \frac{\boldsymbol{\mu}_L \cdot \boldsymbol{\mu}_R}{|\mathbf{R}|^3} - \frac{3(\boldsymbol{\mu}_L \cdot \mathbf{R})(\boldsymbol{\mu}_R \cdot \mathbf{R})}{|\mathbf{R}|^5} \right) .
\end{equation}

The time-dependent effective permittivity $\varepsilon(t)$ is commonly applied during the process. The aqueous environment exhibits a response characterised by a Debye relaxation time $\tau_D \approx 8.3$\,ps. Because this timescale is orders of magnitude longer than the sub-picosecond timescales of excitation transfer, the solvent cannot undergo nuclear equilibration. Thus, the time-dependent permittivity is initially governed by the optical dielectric constant ($\varepsilon_{\rm opt} \approx 1.78$), and the screening effect approaches the static limit ($\varepsilon_{\rm eq} \approx 78$) as given by
\begin{equation}
    \frac{1}{\varepsilon(t)} = \frac{1}{\varepsilon_{\rm eq}} + \left(\frac{1}{\varepsilon_{\rm opt}} - \frac{1}{\varepsilon_{\rm eq}}\right)e^{-t/\tau_D}.
\end{equation}

Assuming a structurally identical homodimer where the ground to first excited state excitation energies are identical for each monomer ($E_L = E_R = E$), the excited state Hamiltonian in the site basis is
\begin{equation}
    \hat H_{\text{site}}(t)=\begin{pmatrix} E & J(t) \\ J(t) & E \end{pmatrix}.
\end{equation}
Diagonalising $\hat H_{\text{site}}(t)$ yields the adiabatic (exciton) basis, where eigenstates are the symmetric (bright) and antisymmetric (dark) combinations such as $\ket{\pm} = \frac{1}{\sqrt{2}}(\ket{1} \pm \ket{2})$. Photoexcitation is treated as a vertical transition directly into the delocalised bright state $\ket{+}$, and its energy separation is the Davydov splitting, given by $\Delta E(t) = 2|J(t)|$. Over several picoseconds, the coupling $J(t)$ relaxes as the solvent environment reorganises from the initial optical dielectric response toward the static equilibrium limit, resulting in a gradual damping of the excitonic interaction.

Because the physical interactions between the system and its aqueous environment are local, we define the master equation (ME) with a Hermitian jump operator $\hat{L} = \sqrt{\hbar \gamma / 2} \hat{\sigma}_z$ in the site basis and the pure dephasing rate $\gamma = 1/T_2^*$ ($\hat{\sigma}_z$ for the Pauli-$z$ operator and $\hbar$ for the reduced Planck constant). For our simulations, a dephasing time of $T_2^* = 60$\,fs was used. While the rigid $\beta$-barrel of fluorescent proteins can partially shield chromophores from solvent fluctuations and extend coherence relative to unbound organic dyes \cite{kim2019venusa206}, this conservatively fast $60$\,fs limit was chosen to align with empirical measurements of rapidly decaying electronic quantum coherence in solvated pigment-protein complexes at ambient temperatures \cite{duan2017nature}. Our choice of jump operator effectively ``measures'' the excitation site, driving the state towards the local site basis, which emerges as the decoherence-resistant pointer basis ($\ket{1}$ and $\ket{2}$) since physical interactions are local. The system's ensemble-averaged density matrix $\hat{\rho}(t)$ with $\ket{1}$ and $\ket{2}$ evolves via the Lindblad equation \cite{manzano2020short}, given by
\begin{equation}
    \frac{d\hat{\rho}}{dt}=-\frac{i}{\hbar}[\hat H_{\text{site}}(t),\hat\rho] + \hat{L} \hat\rho \hat{L} - \frac{1}{2} \{\hat{L}^2, \hat\rho\}\,.
    \label{eq:ME01}
\end{equation}

Building on the Lindblad model, we extend the analysis by simulating individual pure-state trajectories in addition to the ensemble density operator \cite{abrahams2025excitonic}. While the ME provides a deterministic description of the average behaviour of a large ensemble of dimers, it cannot resolve the discrete, incoherent events experienced by a single system. To capture these dynamics, we use a stochastic Schrödinger equation (SSE), which describes the evolution of a single pure-state trajectory $|\psi\rangle$ subject to the continuous environmental monitoring. The density operator is recovered by taking the stochastic average $\mathbb{E}$ over many trajectories, $\hat{\rho} = \mathbb{E}[|\psi\rangle\langle\psi|]$. Then, the SSE is given by
\begin{multline}
    |d\psi\rangle = \left( -\frac{i}{\hbar}\hat{H}_{\text{site}} - \frac{1}{2}(\hat{L} - \langle \hat{L} \rangle)^2 \right) |\psi\rangle dt \\
    + (\hat{L} - \langle \hat{L} \rangle) |\psi\rangle dW \, ,
\end{multline}
where $dW$ is an It\^{o} Wiener process~\cite{petruccione}.

We visualise the individual pure-state trajectories (\textbf{Figures~\ref{fig:ExcitonDynamics}(b) and (c)}) by mapping the system to a Bloch vector $\mathbf{r} = (u, v, w)$ (\textbf{Figure~\ref{fig:ExcitonDynamics}(a)}), where $u = 2\text{Re}(\rho_{12})$, $v = 2\text{Im}(\rho_{12})$, and $w = \rho_{11} - \rho_{22}$. All the numerical simulations are initialised in the delocalised bright state $\ket{+}$, corresponding to a vector on the equator of the Bloch sphere. Individual trajectories exhibit incoherent hopping between sites in the bright-dark basis. These stochastic paths represent pure-state evolution restricted to the surface of the sphere as shown in \textbf{Figure~\ref{fig:ExcitonDynamics}(a)}. Conversely, the ensemble-averaged dynamics governed by the ME in Eq.~(\ref{eq:ME01}) describe the rapid, smooth decay of off-diagonal coherence, leading to a statistical mixture over both sites and evolution toward the centre of the Bloch sphere as the blue solid straight line in \textbf{Figure~\ref{fig:ExcitonDynamics}(a)}.

\section{TDDFT Results}

\begin{table}[t]
\centering
\caption{TDDFT excited states for a monomer used in the coupling analysis. The wavelengths indicate the photon wavelengths corresponding to the transition energies computed in Terachem.}
\label{tab:roots}
\begin{tabular}{cccc}
\toprule
State & Wavelength (nm) & $f$ & $|\mu|$ (a.u.) \\
\midrule
1  & 542.4 & 0.7037 & 3.5449 \\
2  & 476.9 & 0.0238 & 0.6113 \\
3  & 460.2 & 0.0030 & 0.2132 \\
4  & 443.9 & 0.0448 & 0.8091 \\
5  & 391.3 & 0.0304 & 0.6258 \\
6  & 388.8 & 0.0029 & 0.1927 \\
7  & 380.9 & 0.0310 & 0.6235 \\
8  & 340.0 & 0.0304 & 0.5833 \\
9  & 334.7 & 0.0000 & 0.0000 \\
10 & 329.9 & 0.0033 & 0.1893 \\
\bottomrule
\end{tabular}
\end{table}

\begin{figure*}[t]
  \centering
  \begin{subfigure}[b]{0.48\textwidth}
    \centering
    \includegraphics[width=\linewidth]{dimer_view_old_Difference_current.png}
    \caption{}
    \label{fig:density_dimer}
  \end{subfigure}\hfill
  \begin{subfigure}[b]{0.48\textwidth}
    \centering
    \includegraphics[width=\linewidth]{QMRegion.png}
    \caption{}
    \label{fig:density_qm}
  \end{subfigure}
  \caption{Visualisation of the electron transition density $\rho^{g\to e}_m(\mathbf r)$ in Eq.~(\ref{eq:transition_density})  between the ground and first excited state in Table \ref{tab:roots}. The red and blue isosurfaces represent the positive and negative phases of the transition density, mapped onto the dimer frame using PyMOL structural alignment. (a) The full dimer system partitioned into QM and MM domains. The MM embedding environment, treated as point charges, is rendered as a protein backbone with a rainbow colourmap to give a sense of depth perception. (b) A detailed view of the QM region (indicated with sticks and spheres) on the right monomer. This subset is defined by the chromophore, explicit nearest-neighbour residues, and water molecules, allowing for the treatment of electronic polarisation and deprotonation effects. PyMOL \cite{delano2002pymol} generation scripts are available on GitHub.}
  \label{fig:density_combined}
\end{figure*}

Ten time-dependent density functional theory (TDDFT) roots (as eigenenergies) were computed using TeraChem, as shown in Table \ref{tab:roots}. Root 1 was identified as the target excitonic state, located at $542.4~\mathrm{nm}$ with a high oscillator strength of $f=0.7037$ and a transition dipole moment of $|\mu|=3.5449~\mathrm{a.u.}$. To physically characterise this state, the electron transition density ($\rho^{g\to e}$) and the ground-to-excited state electron density difference ($\rho^{e} - \rho^{g}$) were mapped onto the dimer frame (\textbf{Figure~\ref{fig:density_combined}} and \textbf{Figure~\ref{fig:electron_combined}}, respectively). 

The electron transition density in \textbf{Figure~\ref{fig:density_combined}}, which serves as the fundamental quantity for evaluating the Coulombic integral in Eq.~(\ref{eq:J_integral}), exhibits a complex multipolar spatial distribution in 3D that extends across the entire chromophore and adjacent residues. Note that the red and blue isosurfaces in \textbf{Figure~\ref{fig:density_combined}(b)} do not represent static charge gain or loss, but rather the spatial ``lobes'' of the electron transition probability. The spatial overlap of these lobes between the two chromophores determines the magnitude of $J_{\text{TDC}}$, and captures interactions that are ignored when the density is treated as a single point-dipole vector. Concurrently, the electron density difference in \textbf{Figure~\ref{fig:electron_combined}(b)} illustrates the charge reorganisation from the ground to the first excited state (state 1), with the deprotonated phenolic oxygen on the chromophore residue (CR2, anionic form) transferring charge to the rest of the chromophore. 

Then, the electron transition densities were mapped to the dimer frame (centroid separation $27.6~\mathrm{\AA}$). The inter-dipole angle was calculated as $92.85^{\circ}$. The optical dielectric coupling at $t=0$ results are summarised as 

\begin{itemize}
    \item \textbf{Transition Density Coupling:} $J_{\text{TDC}} = 32.82 \pm 1.55~\mathrm{cm^{-1}}$ ($4.07~\mathrm{meV}$),
    \item \textbf{Point-Dipole Approximation:} $J_{\text{PDA}} = 27.64 \pm 1.28~\mathrm{cm^{-1}}$ ($3.43~\mathrm{meV}$),
\end{itemize}

\begin{figure*}[t]
  \centering
  \begin{subfigure}[b]{0.48\textwidth}
    \centering
    \includegraphics[width=\linewidth]{dimer_view_old_Electron_current.png}
    \caption{}
    \label{fig:electron_dimer}
  \end{subfigure}\hfill
  \begin{subfigure}[b]{0.48\textwidth}
    \centering
    \includegraphics[width=\linewidth]{QMRegionElectron.png}
    \caption{}
    \label{fig:electron_qm}
  \end{subfigure}
  \caption{Visualisation of the electron density difference ($\rho^{e} - \rho^{g}$) between the ground and first excited state delocalised over the dimeric Venus geometry from Table \ref{tab:roots}. The green and violet isosurfaces represent electron gain and loss respectively. (a) The dimer frame illustrating the QM/MM partitioning, with the MM backbone rendered in the rainbow colourmap to give a sense of depth perception. (b) A detailed subset of the QM region on the right monomer, highlighting the chromophore and explicit nearest-neighbour residues. The detailed PyMOL \cite{delano2002pymol} codes are available on GitHub.}
  \label{fig:electron_combined}
\end{figure*}

Evaluated on the same 1000 frames, with the same transition density and the same
geometry, the TDC and far-field point-dipole estimates are close. At a centroid
separation of $24.69 \pm 0.32~\mathrm{\AA}$ we obtain
$J_{\mathrm{TDC}} = 32.82 \pm 1.55~\mathrm{cm^{-1}}$ against
$J_{\mathrm{PDA}} = 27.64 \pm 1.28~\mathrm{cm^{-1}}$, a ratio of
$1.187 \pm 0.011$. The spatially extended transition density in
\textbf{Figure~\ref{fig:density_combined}} does allow phase lobes to interact at
effective distances shorter than the centroid separation, but at this separation
the effect is a $\sim 19\%$ correction rather than a qualitative one. The
practical consequence is a useful negative result: for fluorescent-protein
dimers at biological separations the point-dipole approximation is adequate, and
the expensive transition-density integral is not required to estimate $J$.

The coupling must also be compared in a stated screening convention. Our value
carries optical-limit screening, $\varepsilon_{\mathrm{opt}} = 1.77$; the
two-photon polarisation analysis of Cusick et al.\ applies none, disregarding
dielectric attenuation on the grounds of the nanometre scale of the system
\cite{cusick2026polarization}. Placed in both conventions, our
$32.82~\mathrm{cm^{-1}}$ screened corresponds to $58.09~\mathrm{cm^{-1}}$
unscreened, against their $32$--$40~\mathrm{cm^{-1}}$ unscreened
($18.1$--$22.6~\mathrm{cm^{-1}}$ screened). The screening factor cancels from the
ratio, so we are high by $1.61\times$ in either convention; the apparent
coincidence between $32.82$ and $32$--$40$ compares a screened quantity to an
unscreened one and is not an agreement. Inverting their point-dipole expression
for the transition dipole implied by their own measured coupling, separation and
geometry returns $7.6$--$7.9~\mathrm{D}$, against the $9.8~\mathrm{D}$ of our
spectroscopically normalised density; since $J \propto |\mu|^2$, this single
input accounts for most of the residual discrepancy. The choice of
$\varepsilon$ and the transition dipole are therefore the dominant systematics,
both larger than the near-field correction itself.

Finally, the CD couplet does not adjudicate between these values. Nguyen et
al.\ report an apparent coupling of $131$--$186~\mathrm{cm^{-1}}$ inferred from
the peak-to-peak separation of the couplet in a dVenus tandem dimer
\cite{nguyen2025anomalous}. That separation is not an estimator of the exciton
gap: for a splitting small compared with the linewidth the couplet approaches
the derivative of the lineshape, and its extrema are governed by the linewidth
rather than by the splitting. Simulating the couplet through the published
lineshape, the extrema separation is a non-monotonic and two-valued function of
the underlying gap, so a measured separation maps to more than one candidate gap
and cannot invert. Only the couplet amplitude, suppressed by $2|J|/\Omega$,
carries $J$. The same artefact appears in the two-photon polarisation spectrum,
where features separated by $750~\mathrm{cm^{-1}}$ accompany an excitonic
splitting of only $60$--$90~\mathrm{cm^{-1}}$, and was there handled by fitting
amplitudes rather than separations \cite{cusick2026polarization}.

\section{Computational Methods}

The quantum-classical hybrid workflow is automated via a custom Python orchestration that integrates structural preparation, classical relaxation, QM/MM electronic structure calculations, and GPU-accelerated TDC analysis. The full automated procedure is outlined in \textbf{Algorithm \ref{alg:pipeline}}.

The simulation of the dimeric Venus (PDB: 1MYW \cite{berman2000protein}) proceeded through three primary stages as below.

\begin{enumerate}
    \item \textbf{Structural Preparation and Relaxation:} Geometries for both the isolated monomer and the full interacting dimer were prepared and solvated using PDBFixer \cite{openmm2017} and OpenMM \cite{openmm2017}. Classical energy minimisation of both systems was performed to generate a relaxed AMBER14 \cite{amber14_manual} force field environment. The relaxed monomer geometry was extracted to define the Quantum Mechanical (QM) region for subsequent electronic structure calculations, while the relaxed dimer geometry was retained exclusively as the 3D spatial frame for the later coupling analysis (stage 3). To partition the monomer system for QM/MM embedding it is necessary to sever the protein backbone. To avoid artificially introducing chemical instability, these severed bonds were capped with hydrogen atoms.
    
    \item \textbf{Electronic Structure Calculations:} High-level electronic structure calculations were restricted to the isolated monomer geometry. Electronic transition densities and excited-state roots were computed via TDDFT in TeraChem \cite{terachem2009,seritan2021terachem}, treating the monomer's QM region within the point-charge embedding of its classically relaxed MM environment. This stage utilised the range-separated hybrid functional $\omega$B97X-D3/6-311G** \cite{mardirossian2017} and incorporated a non-equilibrium COSMO Polarisable Continuum Model \cite{klamt1993cosmo} for the bulk aqueous solvent.
    
    \item \textbf{Transition Density Coupling (TDC):} The inter-chromophore coupling $J$ was evaluated by reconstructing the full dimeric interaction from the monomeric data. The computed monomeric transition density grids were duplicated and spatially mapped onto the two corresponding chromophore sites of the classically relaxed dimer geometry using PyMOL \cite{delano2002pymol}. Direct numerical integration was then performed across these aligned fine grids (i.e., $\sim 1.5 \times 10^{6}$ TDDFT grid points). Numerical results were subsequently screened by the optical dielectric constant ($\varepsilon_{\rm opt} \approx 1.78$) to account for the sub-picosecond timescale of exciton formation relative to bulk solvent relaxation.
\end{enumerate}

The integration of these structural, electronic, and numerical tools into a combined workflow is detailed in Algorithm \ref{alg:pipeline}.

\section{Summary}

We have characterised the excitonic coupling and energy transfer dynamics in a Venus dimer using a hybrid QM/MM transition-density approach in an open quantum system model. Over a 1000-frame ensemble the coupling is $J = 32.82 \pm 1.55~\mathrm{cm^{-1}}$ under optical-limit screening, corresponding to $2|J| = 65.6~\mathrm{cm^{-1}}$. By integrating distributed transition densities we find that the standard PDA recovers this to within $1.187 \pm 0.011$, so explicit treatment of near-field multipolar effects is not required at biologically relevant separations --- a result that reverses the conclusion of version 1 of this preprint. The coupling is in turn small against the ensemble site-energy detuning, $\langle|\Delta|\rangle = 576~\mathrm{cm^{-1}}$, placing the dimer in a detuning-dominated regime with localised eigenstates.

Our analysis validates the role of non-equilibrium dielectric screening in protein environments. Because the excitonic interaction timescale is significantly faster than the macroscopic Debye relaxation time of bulk water ($\sim 8.3$\,ps), the excitonic coupling can operate in the optical dielectric limit. This inertial ``freezing'' of the solvent protects the initial coupling from overwhelming equilibrium screening with $\varepsilon_{\rm eq}$. The stochastic trajectories illustrate that this coherence is transient and environmental interactions drive the system into a localised pointer basis on a sub-$100$\,fs timescale. This temporal separation reconciles the robust Davydov splitting imprinted at absorption with rapid localisation and subsequent screening within picosecond timescales. Although Fig.~\ref{fig:ExcitonDynamics} well explains the single-photon excited subspace phenomena in the Venus dimer, expanding the time-dependent simulation of excitonic interaction to include higher energy levels would offer a more comprehensive view of non-equilibrium dynamics in fluorescent proteins.

Further refinements to increase the accuracy of our simulation estimates are possible, though they would incur a significant computational cost. Such refinements include enlarging the QM region and treating the phenolic hydrogen fully quantum mechanically to capture proton tunnelling \cite{chakraborty2008development,bourne2024quantum}. Additionally, future work will extend this framework to incorporate explicit vibronic coupling and evaluate the impact of non-Markovian protein phonons on long-term coherence dynamics.

\section{Data and Code Availability}

The computational workflow was executed using the script \texttt{terachem\_full\_pipeline.py} available on GitHub: \url{https://github.com/rchristie95/TerachemVenusA206}. 

\section{Acknowledgements}
R.C and J.J. acknowledge support from the Institute for Information \& Communications Technology Promotion (IITP) grant funded by the Korea government (MSIP) (No. 2019-000003). Y.K. is grateful for support from the G-LAMP Program of the National Research Foundation of Korea, funded by the Ministry of Education (RS-2024-00445180). 

\bibliography{Bibliography}


\begin{algorithm*}[ht]
\caption{Quantum-Classical Hybrid Excitonic Coupling Pipeline \\(\texttt{terachem\_full\_pipeline.py})}
\label{alg:pipeline}
\begin{algorithmic}[1]
\Procedure{PipelineMain}{}
    \State \Call{Stage1\_ClassicalSetup}{}
    \State \Call{Stage2\_TDDFT}{}
    \State \Call{Stage3\_TDC}{}
\EndProcedure

\vspace{0.2cm}
\Procedure{Stage1\_ClassicalSetup}{}
    \State Initialise system with PDBFixer (add missing atoms/hydrogens)
    \State Load AMBER14 XML forcefield parameters
    \State Solvate both dimer and monomer systems 
    \State Relax both solvated geometries via OpenMM classical minimisation, with equilibrium bulk dielectric constant ($\varepsilon_{\rm eq}$) 
    \State Partition system into QM and MM regions (link-atom boundary)
    \State Evaluate formal charges via CCD CIF parsing
    \State Extract MM point charges, applying short-range electrostatic damping
    \State Export monomer QM region coordinates as .xyz file
\EndProcedure

\vspace{0.2cm}
\Procedure{Stage2\_TDDFT}{}
    \State Load the monomer XYZ geometry and MM point charge embeddings
    \State Configure TeraChem settings ($\omega$B97X-D3/6-311G**, COSMO PCM)
    \State Execute TDDFT energy calculation 
    \State Parse excited state roots and select target state by highest oscillator strength
    \State Generate electron transition and difference density grids (DX format)
\EndProcedure

\vspace{0.2cm}
\Procedure{Stage3\_TDC}{}
    \State Load classically relaxed monomer and dimer PDB structures from stage 1
    \State Load the electron transition density from stage 2
    \State Compute transformation matrices $M_L, M_R$ via PyMOL alignment using the dimer geometry 
    \State Parse transition density grid points and charges from .dx file
    \State Map density to Dimer Site L using $M_L$ and Dimer Site R using $M_R$
    \State \Call{CalculateCouplingGPU}{Site L, Site R}
\EndProcedure

\vspace{0.2cm}
\Procedure{CalculateCouplingGPU}{$\mathbf{r}_L, q_L, \mathbf{r}_R, q_R$}
    \State Initialise Numba CUDA or PyOpenCL context
    \State Initialise $J = 0$ 
    \For{each grid point $i$ in Site A accumulate:}
        \State $J \gets J + \sum_j \frac{q_{L,i} \cdot q_{R,j}}{|\mathbf{r}_{L,i} - \mathbf{r}_{R,j}|}$
    \EndFor
    \State Apply optical dielectric screening factor ($\varepsilon_{\mathrm{opt}}$)
    \State \Return Excitonic Coupling $J$
\EndProcedure
\end{algorithmic}
\end{algorithm*}

\end{document}